\chapter{Campi Vettoriali Conservativi}
\section{Cosa devo fare in questo Esercizio?}
\begin{enumerate}
    \item Verificare che il campo sia conservativo 
    \begin{enumerate}
        \item Verificare condizioni necessarie
        \item Verificare condizioni sufficienti
    \end{enumerate}
    \item Calcolare il potenziale
\end{enumerate}

\section{Metodo di integrazione indefinita successiva}
\subsection{Procedura}
\begin{enumerate}
    \item Verificare le condizioni necessarie
    \begin{itemize} 
        \item{\LARGE $\R^{2}$ : $\frac{\partial F_{1}}{\partial y} = \frac{\partial F_{2}}{\partial x}$}
        \item {\LARGE $\R^{3}: \frac{\partial F_1}{\partial y} = \frac{\partial F_2}{\partial x}, \quad 
        \frac{\partial F_1}{\partial z} = \frac{\partial F_3}{\partial x}, \quad 
        \frac{\partial F_2}{\partial z} = \frac{\partial F_3}{\partial y}$} 
    \end{itemize}
    \item Verificare le condizioni sufficienti
    \begin{itemize}
        \item Le componenti di $F$ sono di classe $C^{1}$ su un dominio D
        \item Il campo è irritazionale in $D$
        \item Il dominio $D$ è semplicemente connesso
    \end{itemize}
    \item Porre :
    {\large
    $$
    \vec{\nabla}U=\vec{F}= 
    \begin{cases}
    \frac{\partial U}{\partial x} = \vec{F_{1}} \\
    \frac{\partial U}{\partial y} = \vec{F_{2}} \\
    \frac{\partial U}{\partial z} = \vec{F_{3}} \quad \text{Se siamo in $\R^3$}  \\
    \end{cases}
    $$
    }
    \item Scegliere una componente del campo e calcolarne l'integrale indefinito rispetto alla sua variabile, aggiungere una funzione incognita che dipenda dalle restanti variabili  :  $$U(x, y, z) = \int F_1(x, y, z) \, dx + f(y, z)$$
    \item Si deriva la $U$ appena trovata  rispetto a $y$ e la si uguaglia a $F_{2}$ , si determinerà la forma di $f(y,z)$, la quale conterrà una nuova funzione incognita $g(z)$ dipendente solo dall'ultima variabile
    \item  Si aggiorna l'espressione di $U$, la si deriva rispetto a $z$ e la si uguaglia a $F_3$, semplificando i termini residui e integrando rispetto a $z$, si ricava infine $g(z)$ a meno di una costante numerica reale $C$.
\end{enumerate}

\subsection{Quando usarlo?}
\begin{itemize}
    \item Quando il campo vettoriale ha componenti formate da molti prodotti tra variabili
    \item Quando le funzioni componenti sono facili da derivare parzialmente
    \item Quando il dominio del campo ha "buchi" o esclusioni vicino all'origine degli assi, poiché questo metodo lavora con integrali indefiniti e non richiede di fissare un punto di partenza numerico
\end{itemize}
\subsection{Pro e Contro}
\subsubsection{Pro}
\begin{itemize}
    \item \textbf{Procedimento meccanico}
    \item \textbf{Semplificazione automatica} 
    
    Se il campo è davvero conservativo, durante i passaggi di derivazione i termini misti devono elidersi. Se non succede, ti accorgi immediatamente di aver fatto un errore di calcolo.
\end{itemize}
\subsubsection{Contro}
\begin{itemize}
    \item \textbf{Procedimento lungo}
    \item Richiede molti passaggi algebrici intermedi di derivazione parziale, basta sbagliare un esponente per trascinarsi l'errore fino alla fine
\end{itemize}

\subsection{Esempio 1}
\begin{gather}
\intertext{Calcolare (se esiste) il potenziale di $\vec{F} = (2x y^2 z^2)\vec{e}_1 + (2y x^2 z^2)\vec{e}_2 + (2z x^2 y^2)\vec{e}_3$}
\intertext{Verifichiamo la condizione necessaria:}
\frac{\partial F_1}{\partial y} = \frac{\partial F_2}{\partial x} = 4x y z^2 \\
\frac{\partial F_1}{\partial z} = \frac{\partial F_3}{\partial x} = 4x y^2 z \\ 
\frac{\partial F_2}{\partial z} = \frac{\partial F_3}{\partial y} = 4x^2 y z \\
\intertext{Essendo che le condizioni necessarie sono soddisfatte passiamo a quelle sufficienti} 
F_1 = 2xy^2z^2 \implies D = \{\mathbb{R}^3\} \\
F_2 = 2yx^2z^2 \implies D = \{\mathbb{R}^3\} \\
F_3 = 2zx^2y^2 \implies D = \{\mathbb{R}^3\} \\
D = \{\R^3\}
\intertext{Ora mettiamo tutto a sistema e calcoliamo U}
\vec{\nabla}U = \vec{F} = 
\begin{cases}
    \frac{\partial U}{\partial x} = 2x y^2 z^2 \\
    \frac{\partial U}{\partial y} = 2y x^2 z^2 \\
    \frac{\partial U}{\partial z} = 2z x^2 y^2
\end{cases}
\intertext{Integriamo la prima funzione}
U = \int 2xy^2z^2 \, dx \implies U = x^2y^2z^2 + f(y,z) \\
\intertext{Ora deriviamo U e troviamo f(y,z)}
2x^2yz^2 + \frac{\partial f(y,z)}{\partial y} = 2yx^2z^2 \implies \frac{\partial f(y,z)}{\partial y} = 0 \implies f(y,z) = g(z)\\
2x^2y^2z + g'(z) = 2zx^2y^2 \implies g'(z) = 0 \\
g(z) = c \quad (\text{con } c \in \mathbb{R})
\intertext{Il risultato finale è: } U(x,y,z) = x^2y^2z^2 + c
\end{gather}



\section{Metodo Integrale di Linea lungo i Segmenti Paralleli agli Assi}
\subsection{Procedura}
\begin{gather}
\intertext{Verifichiamo le condizioni necessarie}
\frac{\partial F_1}{\partial y} = \frac{\partial F_2}{\partial x}, \quad 
    \frac{\partial F_1}{\partial z} = \frac{\partial F_3}{\partial x}, \quad 
    \frac{\partial F_2}{\partial z} = \frac{\partial F_3}{\partial y} \\
\intertext{Si fissa un punto appartenente al dominio e si calcola l'integrale tra il punto fisso e un punto generale}
P_{0}(x_{0},y_{0},z_{0}) \quad P(x,y,z) \\
\intertext{Si parametrizza il percorso tra i 2 punti}
\begin{cases}
    x(t) = x_{0} + t(x_{1} - x_{0}) \\
    y(t) = y_{0} + t(y_{1} - y_{0}) \\
    z(t) = z_{0} + t(z_{1} - z_{0})
\end{cases} \quad t \in [0,1] 
\intertext{Calcoliamo il potenziale via integrale di linea di 2ª specie}
U(x,y,z) = \int_{\gamma} \vec{F}(x(t),y(t),z(t  )) \cdot \vec{t} \, dt
\end{gather}

\subsection{Quando usarlo}
\begin{enumerate}
    \item \textbf{Componenti del campo formate da somme isolate}
    \item \textbf{L'origine degli assi fa parte del dominio}
\end{enumerate}
\subsection{Pro e Contro}
\subsubsection{Pro}
\begin{itemize}
    \item \textbf{Veloce}
    \item \textbf{Meno passaggi algrebrici rischiosi}
\end{itemize}
\subsubsection{Contro}
\begin{itemize}
    \item \textbf{Difficile se l'origine non fa parte del dominio}
\end{itemize}

\subsection{Esempio}
\begin{gather}
\intertext{Calcolare (se esiste) il potenziale di $\vec{F} = (2x y^2 z^2)\vec{e}_1 + (2y x^2 z^2)\vec{e}_2 + (2z x^2 y^2)\vec{e}_3$}
\intertext{Verifichiamo la condizione necessaria:}
\frac{\partial F_1}{\partial y} = \frac{\partial F_2}{\partial x} = 4x y z^2 \\
\frac{\partial F_1}{\partial z} = \frac{\partial F_3}{\partial x} = 4x y^2 z \\ 
\frac{\partial F_2}{\partial z} = \frac{\partial F_3}{\partial y} = 4x^2 y z \\
\intertext{Essendo che le condizioni necessarie sono soddisfatte passiamo a quelle sufficienti} 
F_1 = 2xy^2z^2 \implies D = \{\mathbb{R}^3\} \\
F_2 = 2yx^2z^2 \implies D = \{\mathbb{R}^3\} \\
F_3 = 2zx^2y^2 \implies D = \{\mathbb{R}^3\} \\
D = \{\R^3\}
\intertext{Fissiamo un punto }
P_{0} (0,0,0) \quad P(x,y,z) \\
\intertext{Troviamo una parametrizzazione che giugne i 2 punti}
\begin{cases}
    x(t) = x_{0} + t(x_{1} - x_{0}) = xt \\
    y(t) = y_{0} + t(y_{1} - y_{0}) = yt \\
    z(t) = z_{0} + t(z_{1} - z_{0}) = zt
\end{cases} \quad t \in [0,1] 
\intertext{Calcoliamo l'integrale di linea di 2ª specie}
U = \int_{0}^{1}  2x^2 y^2 z^2 t^5 + 2y^2 x^2 z^2 t^5 + 2z^2 x^2 y^2 t^5 \, dt \\
\int_{0}^{1} 6 x^2 y^2 z^2 t^5 \, dt \\
6 x^2 y^2 z^2 \int_{0}^{1} t^5 \, dt \\
6 x^2 y^2 z^2 \cdot [\frac{t^6}{6}]^1_0 \\
U = 6 x^2 y^2 z^2 \cdot \frac{1}{6} = x^2y^2z^2 + c
\end{gather}

\subsection{Esempio 2}
\begin{gather}
    \intertext{Sia } \gamma \intertext{la curva di equazioni parametriche}
    \begin{cases}
        x(t) = t \\
        y(t) = \frac{2t^5}{1-t^5} \\
        z(t) = \frac{4t^10}{1+t^4}
    \end{cases} \quad t \in [0,1]
    \intertext{determinare (se esiste) il valore di  $\lambda$ tale che}  
    \int_{\gamma} [ (\lambda x + 2xy^2z^3)\cdot\frac{dx}{dt} + (\lambda y + 2x^2yz^3)\cdot\frac{dy}{dt} + (\lambda z + 3x^2y^2z^2)\cdot\frac{dz}{dt} ]\, dt = 0
    \intertext{Controlliamo le condizioni necessarie}
    \frac{\partial F_1}{\partial y} = \frac{\partial F_2}{\partial x}, \quad 
    \frac{\partial F_1}{\partial z} = \frac{\partial F_3}{\partial x}, \quad 
    \frac{\partial F_2}{\partial z} = \frac{\partial F_3}{\partial y} \\
    \intertext{Fissiamo il percorso da percorrere}
    P_0{0,0,0} \quad P_{1}(1,1,2) \intertext{$P_{1}$ ottenuto sostituendo $t=1$ nelle funzioni di parametrizzazione}
    \intertext{Parametrizziamo il percorso}
    \begin{cases}
        x(t)= 0 + t (1 - 0) = t \\
        y(t) = 0 + t(1 - 0) = t \\
        z(t) = 0 + t(2 - 0) = 2t
    \end{cases} \quad t \in [0,1]
    \intertext{Ora calcoliamo l'integrale}
    U = \int_{0}^{1} \left[ (\lambda t + 16t^6) + (\lambda t + 16t^6) + (2\lambda t + 12t^6) \right] \, dt \\
\intertext{Sommiamo i termini simili all'interno delle parentesi prima di integrare:}
U = \int_{0}^{1} (6\lambda t + 56t^6) \, dt \\
\intertext{Calcoliamo la primitiva e valutiamo tra gli estremi 0 e 1:}
U = \left[ 3\lambda t^2 + \frac{56}{7}t^7 \right]_0^1 = 3\lambda + 8 \\
\intertext{Poiché il lavoro lungo una curva chiusa (o tra due punti in cui il potenziale si annulla) deve eguagliare la condizione iniziale del problema, imponiamo l'equazione per trovare $\lambda$:}
3\lambda + 8 = 0 \implies \lambda = -\frac{3}{8}
\end{gather}

\subsection{Esempio 3}
\begin{gather}
    \intertext{Calcolare (se esiste) il potenziale del campo vettoriale $\vec{F}=\frac{2x}{x^2+y^2}\vec{e_1}+\frac{2y}{x^2y^2}\vec{e_2}$}
    \intertext{Verificare le condizioni necessarie}
    \frac{\partial F_{1}}{\partial y} = \frac{\partial F_{2}}{\partial x}
    \intertext{Verificare condizioni sufficienti}
    \vec{F_1}:D=\R^2 \backslash \{0,0\} \\
    \vec{F_1}:D=\R^2 \backslash \{0,0\} \\
    D \text{ non è un dominio semplicemente connesso}
    \intertext{Risolvendo otteniamo}
    U = \ln(x^2+y^2)+c
    \intertext{Essendo che il dominio non è semplicemente conesso abbiamo bisgno di conferme aggiuntive, quindi veriifichiamo che}
    \oint_{\gamma} \vec{F} \cdot \vec{t} \, dl = 0 \\ 
    \gamma = 
    \begin{cases}
        x(t)= \cos(t) \\ 
        y(t)= \sin(t) \\
    \end{cases} \quad t \in [0,2\pi] \\ 
    \int^{2\pi}_0 \cos^2(t) + \sin^2(t) = \int^{2\pi}_0 1 \not = 0 \implies \text{Il campo vettoriale non è conservativo} 
\end{gather}